\chapter{Teoria Hodge'a i metoda Bochnera}
\label{ch:hodge-bochner}
\index{teoria Hodge'a}
\index{metoda Bochnera}

Kohomologia de Rhama z Rozdziału~\ref{ch:formy-de-rham}
rozpoznaje formy zamknięte z dokładnością do form dokładnych.
Metryka pozwala wybrać z każdej klasy szczególnego przedstawiciela:
formę harmoniczną. Krzywizna narzuca ograniczenia na takie formy,
a przez nie na topologię rozmaitości. W tym rozdziale rozpiszemy
rachunek geometryczny i dokładnie oddzielimy go od potrzebnego
twierdzenia analitycznego o operatorach eliptycznych.

Przez cały rozdział $(M^n,g)$ jest gładką, zwartą,
zorientowaną rozmaitością riemannowską \emph{bez brzegu},
a $\nabla$ oznacza koneksję Leviego-Civity.
Nie zakładamy spójności, chyba że jest to powiedziane wprost.
Forma objętości $\mathrm{vol}_g$ jest zgodna z wybraną orientacją.
Występowanie brzegu wymagałoby warunków brzegowych i zmieniłoby
twierdzenie o rozkładzie.

\section{Gwiazda Hodge'a i sprzężenie różniczki zewnętrznej}
\label{sec:hodge-gwiazda}

Gwiazda Hodge'a zamienia pomiar w $k$ niezależnych kierunkach
na pomiar w kierunkach do nich prostopadłych. Metryka mówi,
co znaczy „prostopadły” i jak porównywać wielkości;
orientacja wybiera znak. Na przykład w zorientowanej
$\R^3$ z metryką euklidesową $*\dd x=\dd y\wedge\dd z$:
kierunkowi osi $x$ odpowiada prostopadła, zorientowana
płaszczyzna $yz$. Gwiazda nie różniczkuje formy; wykonuje
taką zamianę osobno w każdym punkcie rozmaitości.

Metryka na $T_pM$ indukuje iloczyn skalarny na
$\Lambda^kT_p^*M$: kliny rosnących układów elementów
ortonormalnej kobazy są ortonormalne. Używamy tu algebry
zewnętrznej z Rozdziału~\ref{ch:algebra-tensorowa} i
dodatkowego wyboru metryki.

\begin{definicja}[Gwiazda Hodge'a]
\label{def:hodge-star}
\index{gwiazda Hodge'a}
Dla $\alpha,\beta\in\Omega^k(M)$ forma
$*\beta\in\Omega^{n-k}(M)$ jest jednoznacznie określona przez
\begin{equation}\label{eq:hodge-star-def}
 \alpha\wedge *\beta=\langle\alpha,\beta\rangle_g
                         \,\mathrm{vol}_g.
\end{equation}
Jeśli $(e^1,\ldots,e^n)$ jest dodatnio zorientowaną
ortonormalną kobazą i $I=(i_1<\cdots<i_k)$, to
\[
 *e^I=\varepsilon(I,I^c)e^{I^c},
\]
gdzie $\varepsilon(I,I^c)$ jest znakiem permutacji
$(i_1,\ldots,i_k,I^c)$.
\end{definicja}

Istnienie wynika z ostatniego wzoru: po rozłożeniu form w
ortonormalnej kobazie współczynnik przy $\mathrm{vol}_g$
po lewej stronie~\eqref{eq:hodge-star-def} jest właśnie
ich iloczynem skalarnym. Wzór nie zależy od wyboru
dodatnio zorientowanej kobazy, bo obie strony definicji
są niezmiennicze. Gładkość $*$ wynika z gładkości metryki.

\begin{stwierdzenie}[Dwa zastosowania gwiazdy]
\label{prop:hodge-star-square}
Na $k$-formach zachodzi
\begin{equation}\label{eq:hodge-star-square}
 **\alpha=(-1)^{k(n-k)}\alpha,
 \qquad *1=\mathrm{vol}_g,\quad *\mathrm{vol}_g=1.
\end{equation}
\end{stwierdzenie}
\begin{proof}
W dodatniej ortonormalnej kobazie pierwsze zastosowanie $*$
przenosi $e^I$ na dopełniający klin ze znakiem permutacji.
Drugie przestawia blok długości $n-k$ za blok długości $k$.
Wymaga to $k(n-k)$ transpozycji, więc iloczyn obu znaków
jest $(-1)^{k(n-k)}$. Dla pustego i pełnego układu indeksów
otrzymujemy dwie pozostałe równości.
\end{proof}

\begin{przyklad}[Płaszczyzna i torus]
W dodatniej ortonormalnej kobazie $(\dd x,\dd y)$ mamy
$*\dd x=\dd y$, $*\dd y=-\dd x$, a
$*(\dd x\wedge\dd y)=1$. Te same formuły obowiązują na
płaskim torusie $\R^2/\Z^2$ dla form pochodzących z
niezmienniczych względem przesunięć form płaszczyzny.
Znak $*\dd y$ zapewnia $**=-1$ na $1$-formach.
\end{przyklad}

Definiujemy iloczyn $L^2$ gładkich $k$-form przez
\begin{equation}\label{eq:hodge-l2}
 (\alpha,\beta)_{L^2}
 =\int_M\langle\alpha,\beta\rangle_g\,\dd\mu_g
 =\int_M\alpha\wedge*\beta.
\end{equation}
Równość z prawej wynika z~\eqref{eq:hodge-star-def}
i zgodności $\dd\mu_g$ z $\mathrm{vol}_g$ opisanej w
Stwierdzeniu~\ref{prop:calka-skierowana-forma}.

\begin{definicja}[Kodyferencjał i laplasjan Hodge'a]
\label{def:hodge-codifferential}
\index{kodyferencjal@Kodyferencjał}
\index{laplasjan Hodge'a}
Na $k$-formach kładziemy
\begin{equation}\label{eq:hodge-delta}
 \delta=(-1)^{n(k+1)+1}*\dd*:\Omega^k(M)\longrightarrow
                                     \Omega^{k-1}(M),
 \qquad\delta|_{\Omega^0}=0.
\end{equation}
Laplasjanem Hodge'a jest
\begin{equation}\label{eq:hodge-laplacian}
 \Delta_H=\dd\delta+\delta\dd.
\end{equation}
Forma $\alpha$ jest \emph{harmoniczna}, gdy
$\Delta_H\alpha=0$; przestrzeń harmonicznych $k$-form
oznaczamy przez $\mathcal H^k(M,g)$.
\end{definicja}

Kodyferencjał $\delta$ mierzy drugą stronę zmienności
formy. Dla $1$-formy $\alpha$ odpowiadającej polu
$\alpha^\sharp$ jest ujemną dywergencją tego pola:
$\delta\alpha=-\operatorname{div}(\alpha^\sharp)$.
Różniczka $\dd$ wykrywa natomiast jego obieg
(w wymiarze trzy odpowiada rotacji). Laplasjan Hodge'a
łączy oba pomiary: $\Delta_H=\dd\delta+\delta\dd$.
Na funkcjach pozostaje ujemnym laplasjanem znanym z
rachunku gradientu. Na zwartej rozmaitości bez brzegu
zerowa energia tego operatora oznacza
$\dd\alpha=\delta\alpha=0$; dla $1$-form jest to
jednoczesny brak obiegu i dywergencji.
Wykażemy to poniżej.

Symbol $\delta$ oznacza tutaj operator na \emph{formach}.
W Rozdziale~\ref{ch:formy-de-rham} ten sam znak występuje
także przy kobrzegu kochain; dziedzina rozstrzyga znaczenie.
Na funkcjach $\Delta_H f=\delta\dd f=-\operatorname{div}
(\nabla_g f)$, czyli znak jest przeciwny do konwencji
$\Delta f=\operatorname{div}(\nabla f)$ z
Rozdziału~\ref{ch:krzywizna}.

\begin{stwierdzenie}[Sprzężenie i energia]
\label{prop:hodge-adjoint-energy}
Dla $1\leq k\leq n$, $\eta\in\Omega^{k-1}(M)$ i
$\alpha\in\Omega^k(M)$ zachodzi
\begin{equation}\label{eq:hodge-adjoint}
 (\dd\eta,\alpha)_{L^2}=(\eta,\delta\alpha)_{L^2}.
\end{equation}
Przy umowie $\Omega^{-1}(M)=\Omega^{n+1}(M)=0$ dla
każdej $\alpha\in\Omega^k(M)$, $0\leq k\leq n$, mamy
\begin{equation}\label{eq:hodge-energy}
 (\Delta_H\alpha,\alpha)_{L^2}
 =\|\dd\alpha\|_{L^2}^2+\|\delta\alpha\|_{L^2}^2,
\end{equation}
więc $\alpha$ jest harmoniczna wtedy i tylko wtedy, gdy
$\dd\alpha=0$ i $\delta\alpha=0$.
\end{stwierdzenie}
\begin{proof}
Twierdzenie Stokesa na $M$ bez brzegu daje
\[
 0=\int_M\dd(\eta\wedge*\alpha)
  =\int_M\dd\eta\wedge*\alpha
   +(-1)^{k-1}\int_M\eta\wedge\dd*\alpha.
\]
Z~\eqref{eq:hodge-star-square} i wykładnika w
\eqref{eq:hodge-delta} wynika
$*\delta\alpha=(-1)^k\dd*\alpha$.
Po podstawieniu druga całka jest
$-(\eta,\delta\alpha)_{L^2}$, co dowodzi
\eqref{eq:hodge-adjoint}. Stosujemy tę równość raz do
$\dd\delta\alpha$, a raz do $\delta\dd\alpha$,
uzyskując~\eqref{eq:hodge-energy}; na krańcach zakresu
pomijamy zerowe składniki.
Jeżeli jej prawa strona znika, obie nieujemne całki są
zerowe. Gładkość form sprawia, że $\dd\alpha$ i
$\delta\alpha$ znikają w każdym punkcie. Implikacja
odwrotna jest bezpośrednia.
\end{proof}

Ponieważ $\dd^2=0$, a z~\eqref{eq:hodge-delta} także
$\delta^2=0$, otrzymujemy $\dd\Delta_H=\Delta_H\dd$
oraz $\delta\Delta_H=\Delta_H\delta$. Równania te
pozwalają przenosić rozkład harmoniczny między stopniami.

\section{Eliptyczność i rozkład Hodge'a}
\label{sec:hodge-rozklad}

Samo równanie energii nie daje istnienia harmonicznego
przedstawiciela klasy: trzeba jeszcze rozwiązać równanie
różniczkowe $\Delta_Hu=f$. Jego lokalny charakter wyjaśnia
symbol główny.

\begin{stwierdzenie}[Symbol laplasjanu]
\label{prop:hodge-symbol}
Jeśli $\xi\in T_p^*M$ jest niezerowy, a w symbolu operatora
pochodne zastępujemy przez $i\xi$, to
\begin{equation}\label{eq:hodge-symbol}
 \sigma_2(\Delta_H)(p,\xi)
 =|\xi|_g^2\,\id_{\Lambda^kT_p^*M}.
\end{equation}
Zatem $\Delta_H$ jest operatorem eliptycznym.
\end{stwierdzenie}
\begin{proof}
W najwyższym rzędzie symbol $\dd$ jest
$i\xi\wedge(\cdot)$, a symbol $\delta$ jest
$-i\iota_{\xi^\sharp}(\cdot)$; drugi wzór można
odczytać z~\eqref{eq:hodge-delta} w ortonormalnej kobazie.
Dlatego symbol sumy $\dd\delta+\delta\dd$ jest
\[
 \xi\wedge\iota_{\xi^\sharp}
 +\iota_{\xi^\sharp}(\xi\wedge\cdot).
\]
Na dowolnej formie $\beta$ relacja
$\iota_{\xi^\sharp}(\xi\wedge\beta)
=|\xi|^2\beta-\xi\wedge\iota_{\xi^\sharp}\beta$
redukuje tę sumę do $|\xi|^2\beta$. Dla $\xi\ne0$
jest to izomorfizm, zgodnie z definicją eliptyczności.
\end{proof}

\begin{twierdzenie}[Analityczne wejście: alternatywa eliptyczna]
\label{thm:hodge-analytic-input}
Dla operatora $\Delta_H$ na $k$-formach zwartej
rozmaitości bez brzegu przestrzeń $\mathcal H^k(M,g)$
jest skończenie wymiarowa i składa się z form gładkich.
Jeżeli gładka forma $f$ jest prostopadła do
$\mathcal H^k(M,g)$ w $L^2$, to istnieje jednoznaczna
gładka forma $u\perp\mathcal H^k(M,g)$ taka, że
$\Delta_Hu=f$.
\end{twierdzenie}

Jest to \emph{przyjęte twierdzenie analityczne},
nie wynik dowiedziony w tym skrypcie. Dowód wymaga
oszacowań eliptycznych, zwartości zanurzeń Sobolewa i
regularności słabych rozwiązań. Rozkład Hodge'a i rolę
parametryksu omawiają
\href{https://ocw.mit.edu/courses/18-966-geometry-of-manifolds-spring-2007/d6848bb391c032ef27993e984fef4558_lect15.pdf}
{notatki MIT, wykład 15}; klasyczne ujęcie pochodzi z pracy
\href{https://doi.org/10.2307/1969552}
{K.\ Kodairy, \emph{Harmonic Fields in Riemannian Manifolds}
(1949)}. Wszystkie poniższe wnioski z tego wejścia
udowodnimy już w ramach rachunku form.

\begin{twierdzenie}[Rozkład Hodge'a]
\label{thm:hodge-decomposition}
Każda $\omega\in\Omega^k(M)$ ma jednoznaczny rozkład
ortogonalny w $L^2$
\begin{equation}\label{eq:hodge-decomposition}
 \omega=h+\dd\eta+\delta\zeta,
 \qquad h\in\mathcal H^k(M,g),\quad
 \eta\in\Omega^{k-1}(M),\quad
 \zeta\in\Omega^{k+1}(M).
\end{equation}
Jednoznaczne są trzy \emph{składniki} $h$, $\dd\eta$,
$\delta\zeta$; same potencjały $\eta,\zeta$
nie muszą być jednoznaczne.
\end{twierdzenie}
\begin{proof}
Rzut ortogonalny $h$ na skończenie wymiarową
$\mathcal H^k$ jest gładki. Z
Twierdzenia~\ref{thm:hodge-analytic-input} istnieje
gładkie $u$ z $\Delta_Hu=\omega-h$.
Stąd $\omega=h+\dd(\delta u)+\delta(\dd u)$.
Harmoniczna forma jest zamknięta i współzamknięta na mocy
\eqref{eq:hodge-energy}, więc z~\eqref{eq:hodge-adjoint}
jest prostopadła do form dokładnych i współdokładnych.
Te dwie ostatnie przestrzenie też są ortogonalne:
\[
 (\dd\eta,\delta\zeta)_{L^2}
 =(\dd^2\eta,\zeta)_{L^2}=0.
\]
Jeśli dwie trójki składników dawałyby tę samą formę,
ich różnica byłaby sumą trzech wzajemnie prostopadłych
form o sumie zerowej. Kwadrat normy tej sumy jest sumą
kwadratów ich norm; wszystkie muszą zniknąć.
\end{proof}

\begin{wniosek}[Harmoniczny przedstawiciel klasy]
\label{cor:hodge-de-rham}
Odwzorowanie
\[
 \mathcal H^k(M,g)\longrightarrow H^k_{\mathrm{dR}}(M),
 \qquad h\longmapsto[h]
\]
jest izomorfizmem. Każda klasa ma więc dokładnie jeden
harmoniczny przedstawiciel, a liczba Bettiego
$b_k=\dim H^k_{\mathrm{dR}}(M)$ jest skończona.
\end{wniosek}
\begin{proof}
Jeśli $\omega$ jest zamknięta, bierzemy jej rozkład
\eqref{eq:hodge-decomposition}. Po sparowaniu z
$\delta\zeta$ dostajemy
\[
 \|\delta\zeta\|_{L^2}^2
 =(\omega,\delta\zeta)_{L^2}
 =(\dd\omega,\zeta)_{L^2}=0,
\]
bo pozostałe składniki są do $\delta\zeta$
ortogonalne. Zatem $\omega=h+\dd\eta$ i $[\omega]=[h]$.
Jeżeli harmoniczna $h$ jest dokładna, powiedzmy
$h=\dd\eta$, to
$\|h\|_{L^2}^2=(\delta h,\eta)_{L^2}=0$.
To daje surjektywność i injektywność. Skończoność
wynika z Twierdzenia~\ref{thm:hodge-analytic-input}.
\end{proof}

Wybór metryki zmienia \emph{formę} harmonicznego
przedstawiciela, ale nie grupę kohomologii.
Geometrycznie $h$ jest przedstawicielem o najmniejszej
normie $L^2$ w swojej klasie. Rzeczywiście każdy inny
przedstawiciel ma postać $h+\dd\eta$, a z
$\delta h=0$ i~\eqref{eq:hodge-adjoint} wynika
\[
 \|h+\dd\eta\|_{L^2}^2
 =\|h\|_{L^2}^2+\|\dd\eta\|_{L^2}^2.
\]
Na torusie taka forma może być niezerowa: $\dd x$
nie ma lokalnego obiegu ani dywergencji, lecz jej całka
wzdłuż pętli okrążającej torus jest niezerowa.
Gwiazda $*$ jest izometrią punktową i, przez
\eqref{eq:hodge-delta}, przenosi formy harmoniczne
stopnia $k$ na formy harmoniczne stopnia $n-k$.
Otrzymujemy między innymi $b_k=b_{n-k}$.

\section{Wzór Weitzenböcka dla form stopnia pierwszego}
\label{sec:bochner-weitzenbock}

Laplasjan Hodge'a używa różniczki zewnętrznej.
Koneksja daje inny operator: laplasjan szorstki
$\nabla^*\nabla=-\operatorname{tr}_g\nabla^2$, gdzie
$\nabla^2_{X,Y}\alpha
=\nabla_X\nabla_Y\alpha-\nabla_{\nabla_XY}\alpha$.
Różnicę między nimi mierzy krzywizna Ricciego
z Definicji~\ref{def:ricci-skalarna}.

\begin{twierdzenie}[Wzór Weitzenböcka na $1$-formach]
\label{thm:weitzenbock-one-forms}
Przy konwencji krzywizny~\eqref{eq:riemann-def},
dla każdej $\alpha\in\Omega^1(M)$ mamy
\begin{equation}\label{eq:weitzenbock-one-forms}
 \Delta_H\alpha=\nabla^*\nabla\alpha
                      +\mathrm{Ric}^{\sharp}(\alpha),
 \qquad
 (\mathrm{Ric}^{\sharp}\alpha)_j
   =\mathrm{Ric}_j{}^a\alpha_a.
\end{equation}
\end{twierdzenie}
\begin{proof}
W punkcie $p$ wybieramy lokalną ortonormalną ramkę
$(e_i)$ z $\nabla e_i(p)=0$. Jeśli
$\alpha_i=\alpha(e_i)$, to w tym punkcie
\[
 (\dd\alpha)_{ij}=\nabla_i\alpha_j-\nabla_j\alpha_i,
 \qquad \delta\alpha=-\nabla^i\alpha_i,
 \qquad (\delta\beta)_j=-\nabla^i\beta_{ij}.
\]
Ostatni wzór dla $2$-formy $\beta$ wynika z
\eqref{eq:hodge-delta} w tej ramce. Wstawiając go do
$\dd\delta\alpha+\delta\dd\alpha$, dostajemy
\begin{align*}
 (\Delta_H\alpha)_j
 &=-\nabla_j\nabla^i\alpha_i
   -\nabla^i(\nabla_i\alpha_j-\nabla_j\alpha_i)\\
 &=-\nabla^i\nabla_i\alpha_j
   +(\nabla^i\nabla_j-\nabla_j\nabla^i)\alpha_i.
\end{align*}
Koneksja na kowektorach ma znak przeciwny do koneksji
na wektorach, toteż
\[
 ([\nabla_i,\nabla_j]\alpha)_k
 =-R^a{}_{kij}\alpha_a.
\]
Po zwężeniu pierwszego i trzeciego indeksu używamy
antysymetrii tensora krzywizny w dwóch ostatnich
argumentach $\mathcal R$ oraz~\eqref{eq:ricci-def}:
\[
 -\sum_i R^a{}_{iij}\alpha_a
 =\sum_iR^i{}_{aij}\alpha_a
 =\mathrm{Ric}_j{}^a\alpha_a.
\]
Pierwszy składnik wcześniejszego rachunku to w ramce
normalnej $(\nabla^*\nabla\alpha)_j$. Ponieważ oba
człony są tensorowe, równość zachodzi na całym $M$.
\end{proof}

\begin{stwierdzenie}[Zintegrowana tożsamość Bochnera]
\label{prop:bochner-integrated}
Dla każdej $1$-formy $\alpha$ zachodzi
\begin{equation}\label{eq:bochner-integrated}
 \|\dd\alpha\|_{L^2}^2+\|\delta\alpha\|_{L^2}^2
 =\|\nabla\alpha\|_{L^2}^2
  +\int_M\mathrm{Ric}(\alpha^\sharp,\alpha^\sharp)
                       \,\dd\mu_g.
\end{equation}
\end{stwierdzenie}
\begin{proof}
Parujemy~\eqref{eq:weitzenbock-one-forms} z $\alpha$
i całkujemy. Lewa strona jest lewą stroną
\eqref{eq:bochner-integrated} dzięki
\eqref{eq:hodge-energy}. Dla członu z koneksją
stosujemy twierdzenie o dywergencji z
Rozdziału~\ref{ch:metryka-riemannowska} do pola
określonego przez $g(X,Y)=\langle\nabla_Y\alpha,\alpha\rangle$.
W ramce normalnej
\[
 \operatorname{div}X
 =\langle\operatorname{tr}_g\nabla^2\alpha,\alpha\rangle
  +|\nabla\alpha|^2.
\]
Całka dywergencji znika, ponieważ $M$ nie ma brzegu.
Pozostały składnik to dokładnie
$\mathrm{Ric}(\alpha^\sharp,\alpha^\sharp)$.
\end{proof}

\begin{wniosek}[Znikanie Bochnera]
\label{cor:bochner-vanishing}
Jeżeli $\mathrm{Ric}\geq0$, każda harmoniczna $1$-forma
jest równoległa. Jeśli ponadto $M$ jest spójna i
$\mathrm{Ric}$ jest dodatnio określony choćby w jednym
punkcie, to $H^1_{\mathrm{dR}}(M)=0$.
W szczególności wynika to z warunku $\mathrm{Ric}>0$
na całym $M$.
\end{wniosek}
\begin{proof}
Harmoniczność zeruje lewą stronę
\eqref{eq:bochner-integrated}. Przy $\mathrm{Ric}\geq0$
obie całki po prawej są nieujemne, zatem
$\nabla\alpha=0$. Długość formy równoległej jest
stała na każdej składowej spójnej. Jeśli Ricci jest
dodatnio określony w punkcie składowej, równość
$\mathrm{Ric}(\alpha^\sharp,\alpha^\sharp)=0$ tamże
wymusza $\alpha=0$ w tym punkcie, a więc wszędzie
na składowej. Na koniec używamy
Wniosku~\ref{cor:hodge-de-rham}.
\end{proof}

\begin{przyklad}[Sfera i płaski torus]
Na okrągłej $S_a^n$, $n\geq2$, Rozdział~\ref{ch:krzywizna}
daje $\mathrm{Ric}=(n-1)a^{-2}g$.
Z Wniosku~\ref{cor:bochner-vanishing} dostajemy
$H^1_{\mathrm{dR}}(S_a^n)=0$; jest to analityczny
dowód topologicznego wyniku obliczonego wcześniej.
Na płaskim $T^n=\R^n/\Z^n$ tensor Ricciego jest
zerowy. Formy $\dd x^1,\ldots,\dd x^n$ schodzą na torus
i są równoległe, więc harmoniczne. Każda harmoniczna
$1$-forma jest równoległa, a jej współczynniki w tej
globalnej kobazie są stałe. Stąd
$\dim H^1_{\mathrm{dR}}(T^n)=n$.
W szczególności torus pokazuje, dlaczego sam warunek
$\mathrm{Ric}\geq0$ nie wystarcza do znikania $H^1$.
\end{przyklad}

\section{Cztery wymiary: formy samodualne i sygnatura}
\label{sec:hodge-four-manifolds}

W wymiarze cztery~\eqref{eq:hodge-star-square} daje
$*^2=1$ na $2$-formach. Stąd każda $2$-forma rozkłada się
punktowo na
\[
 \alpha^+=\tfrac12(\alpha+*\alpha),\qquad
 \alpha^-=\tfrac12(\alpha-*\alpha),\qquad
 *\alpha^\pm=\pm\alpha^\pm.
\]
Obie podwiązki $\Lambda^2_+$ i $\Lambda^2_-$ mają rangę
trzy. Na przykład w dodatniej ortonormalnej kobazie
$e^{12}+e^{34}$ jest samodualna, a $e^{12}-e^{34}$
jest antysamodualna. Jest to algebraiczna podstawa
późniejszych zagadnień czterowymiarowych.
Dla $M^4$ przez $\operatorname{sign}(M)$ oznaczamy sygnaturę formy
przecięcia $Q_M\otimes\R$ z
Rozdziału~\ref{ch:przeciecia-dualnosc-klasy}: liczbę
dodatnich minus liczbę ujemnych kwadratów po jej
diagonalizacji nad $\R$.

\begin{stwierdzenie}[Harmoniczny rozkład formy przecięcia]
\label{prop:hodge-signature}
Jeżeli $M^4$ jest spójna, to $*$ zachowuje
$\mathcal H^2(M,g)$ i daje rozkład
$\mathcal H^2=\mathcal H^2_+\oplus\mathcal H^2_-$.
Rzeczywiste rozszerzenie formy przecięcia z
Rozdziału~\ref{ch:przeciecia-dualnosc-klasy}
ma na tych przestrzeniach odpowiednio znak dodatni
i ujemny. W konsekwencji
\begin{equation}\label{eq:hodge-signature}
 \operatorname{sign}(M)
 =\dim\mathcal H^2_+-\dim\mathcal H^2_-.
\end{equation}
\end{stwierdzenie}
\begin{proof}
Na $2$-formach w wymiarze cztery
$\delta=-*\dd*$. Jeżeli $\alpha$ jest harmoniczna,
to $\dd\alpha=\delta\alpha=0$, więc
$\dd*\alpha=0$ oraz $\delta*\alpha=-*\dd\alpha=0$.
Zatem $*$ działa na harmonicznych $2$-formach, a
ponieważ $*^2=1$, dzieli je na dwa podane podprzestrzenie.
Twierdzenie de Rhama i dualność Poincarégo
utożsamiają rozszerzenie całkowitej formy
$Q_M$ do $\R$ z parowaniem
$([\alpha],[\beta])\mapsto\int_M\alpha\wedge\beta$
na $H^2_{\mathrm{dR}}(M)$. Wniosek~\ref{cor:hodge-de-rham}
pozwala liczyć je na harmonicznych przedstawicielach:
\[
 (Q_M\otimes\R)([\alpha],[\beta])=\int_M\alpha\wedge\beta
 =\int_M\langle\alpha,*\beta\rangle\,\dd\mu_g.
\]
Na $\mathcal H^2_+$ jest to iloczyn $L^2$, na
$\mathcal H^2_-$ jego przeciwieństwo, a składniki
różnych znaków są ortogonalne. Rozszerzenie do $\R$
zachowuje sygnaturę formy całkowitej, co
daje~\eqref{eq:hodge-signature}.
\end{proof}

\begin{przyklad}[Dwie sfery]
Na $S^2\times S^2$ z metryką produktową i orientacją
produktową niech $\omega_1,\omega_2$ będą cofnięciami
form objętości czynników. Wtedy $*\omega_1=\omega_2$
i $*\omega_2=\omega_1$. Obie formy są zamknięte,
a przez te równości także współzamknięte, zatem harmoniczne.
Formy $\omega_1+\omega_2$ i $\omega_1-\omega_2$ należą
odpowiednio do $\mathcal H^2_+$ i $\mathcal H^2_-$.
Ponieważ $b_2(S^2\times S^2)=2$, każda z tych przestrzeni
ma wymiar jeden i $\operatorname{sign}(S^2\times S^2)=0$,
zgodnie z obliczeniem formy przecięcia w
Rozdziale~\ref{ch:przeciecia-dualnosc-klasy}.
\end{przyklad}

\section{Trzy wzory indeksowe potrzebne dalej}
\label{sec:hodge-indeks}
\index{indeks operatora eliptycznego}
\index{twierdzenie Atiyaha--Singera}

W dalszych rozdziałach potrzebujemy trzech konkretnych
wyników: wzoru na sygnaturę w wymiarach cztery i osiem,
całkowitości liczby \(\widehat A\) dla zamkniętej spinowej
\(8\)-rozmaitości oraz indeksu operatora deformacji ASD.
Udowodnimy je wspólnym rachunkiem jądra cieplnego.
Nie dowodzimy tu twierdzenia indeksowego dla dowolnego
symbolu eliptycznego. Jedynym przyjętym wejściem
analitycznym jest standardowy rachunek spektralny
operatorów eliptycznych i oszacowanie gaussowskie
jądra cieplnego, użyte jawnie w poniższym dowodzie.

\subsection{Indeks i superślad cieplny}
\label{subsec:hodge-indeks-cieplo}

Jeśli \(D:C^\infty(E^+)\to C^\infty(E^-)\) jest operatorem
eliptycznym między wiązkami tej samej rangi nad
zamkniętą rozmaitością, jego jądro i kojądro są
skończenie wymiarowe. Liczbę
\[
 \operatorname{ind}D
 =\dim_{\C}\ker D-\dim_{\C}\ker D^*
\]
nazywamy \emph{indeksem}; dla wiązek rzeczywistych
liczymy wymiary nad \(\R\). Skończoność i dyskretność
widma niżej używanych operatorów są częścią
przyjętego rachunku eliptycznego. Symbol pierwszego
rzędu powstaje przez zastąpienie \(\partial_j\) przez
\(i\xi_j\), tak jak przed
Stwierdzeniem~\ref{prop:hodge-symbol}. W szczególności
symbol operatora Diraca jest \(i\,c(\xi)\); relacja
\(c(\xi)^2=-|\xi|^2\) daje jego odwracalność dla
\(\xi\ne0\).

\begin{lemat}[Tożsamość McKeana--Singera]
\label{lem:hodge-mckean-singer}
Niech \(D:E^+\to E^-\) będzie operatorem eliptycznym
pierwszego rzędu na zamkniętej rozmaitości, a \(K_t^\pm\)
jądrami cieplnymi odpowiednio \(D^*D\) i \(DD^*\).
Wtedy dla każdego \(t>0\)
\begin{equation}\label{eq:hodge-mckean-singer}
 \operatorname{ind}D
 =\operatorname{Tr}e^{-tD^*D}
  -\operatorname{Tr}e^{-tDD^*}
 =\int_M\operatorname{str}K_t(x,x)\,d\mu_g(x).
\end{equation}
\end{lemat}
\begin{proof}
Jeśli \(D^*Du=\lambda u\) i \(\lambda>0\), to
\(Du\ne0\) oraz \(DD^*(Du)=\lambda Du\).
Odwrotny izomorfizm na przestrzeniach własnych
daje \(\lambda^{-1}D^*\). Ich wkłady do różnicy
śladów znoszą się. Dla \(\lambda=0\) pozostają
\(\ker D\) i \(\ker D^*\). Drugą równość otrzymujemy,
pisząc ślad operatora o gładkim jądrze jako całkę
jego śladu na przekątnej. Zwarta rozmaitość
zapewnia zbieżność śladów dla \(t>0\).
\end{proof}

Przypomnijmy z Rozdziału~\ref{ch:grass-spin}
spinory chiralne \(S=S^+\oplus S^-\).
Na rozmaitości spinowej, dla hermitowskiej wiązki
\(W\) z koneksją, operator
\[
 D_W^+=c\circ\nabla^{S^+\otimes W}:
 \Gamma(S^+\otimes W)\longrightarrow
 \Gamma(S^-\otimes W)
\]
jest skręconym operatorem Diraca. Na małej kuli
takie spinory i tak istnieją; dlatego ten sam
rachunek lokalny ma sens również wtedy, gdy
globalny operator pochodzi z modułu Clifforda
na rozmaitości niespinowej.

\begin{lemat}[Lokalny rachunek jądra cieplnego]
\label{lem:hodge-local-index}
Niech \(M^{2m}\) będzie zamkniętą rozmaitością
spinową, a \(W\) hermitowską wiązką z koneksją.
Dla jądra cieplnego operatora Diraca \(D_W\)
zachodzi punktowo, w normalizacji Cherna--Weila,
\begin{equation}\label{eq:hodge-local-index}
 \lim_{t\downarrow0}\operatorname{str}K_t(x,x)\,d\mu_g
 =\bigl[\widehat A(TM,\nabla^{TM})\,
        \operatorname{ch}(W,\nabla^W)\bigr]_{2m}.
\end{equation}
To samo lokalne obliczenie stosuje się do
globalnego modułu Clifforda, który na kulach
ma postać \(S\otimes W\), nawet gdy \(S\) i \(W\)
osobno nie sklejają się globalnie.
\end{lemat}
\begin{proof}
Opisujemy także krok analityczny odpowiedzialny za
przejście od lokalnego modelu do granicy.
Ustalmy punkt \(x\), normalne współrzędne \(y\)
i ramy otrzymane przez transport równoległy
po promieniach. Metryka ma postać
\(g_{ij}(y)=\delta_{ij}+O(|y|^2)\).
W ramie radialnej forma koneksji wiązki \(W\)
spełnia
\[
 A_i(y)=-\tfrac12F^W_{ij}(x)y^j+O(|y|^2).
\]
Wynika to z warunku \(y^iA_i(y)=0\): pochodna
wzdłuż promienia daje krzywiznę, a całkowanie
jej wartości liniowej daje czynnik \(1/2\).
Analogiczne rozwinięcie koneksji spinowej
ma w symbolach postać
\(\nabla_i^S=\partial_i-\frac14\Omega_{ij}(x)y^j
+O(|y|^2)\), gdzie
\(\Omega_{ij}=\frac12\sum_{a,b}R_{ijab}e^a\wedge e^b\).

Wzór Weitzenböcka dla skręconego Diraca wyraża
\(D_W^2\) jako laplasjan koneksji, składnik
\(\operatorname{scal}/4\) i działanie krzywizny \(W\)
przez dwa mnożenia Clifforda. Można go sprawdzić
w ramie normalnej: w kwadracie
\(\sum_i c(e_i)\nabla_i\) część symetryczna daje
laplasjan, a komutatory
\([\nabla_i,\nabla_j]\) dają krzywizny.
Utożsammy teraz \(\operatorname{End}(S_x)\)
z \(\Lambda^*T_x^*M\otimes\C\) za pomocą
symbolu Clifforda. Iloczyn symboli \(a\) i \(b\)
ma część najwyższego stopnia \(a\wedge b\);
pozostałe części mają stopień niższy co najmniej
o dwa.

Podstawmy \(y=\varepsilon z\),
\(t=\varepsilon^2s\) i nadajmy formie stopnia
\(r\) wagę \(\varepsilon^{-r}\).
Dokładniej, na jednomianie \(z^\alpha a_r\)
operator \(\mathcal T_\varepsilon\) mnoży przez
\(\varepsilon^{|\alpha|-r}\), a
\(H_\varepsilon=\varepsilon^2\mathcal T_\varepsilon
D_W^2\mathcal T_\varepsilon^{-1}\).
W tym
\emph{skalowaniu Getzlera} działanie generatora
Clifforda na symbolu \(a\) ma postać
\(c_\varepsilon(v)a=v\wedge a
-\varepsilon^2\iota_v a\).
Kontrakcje znikają w granicy, a krzywizny zostają.
Przeskalowany kwadrat Diraca zbiega na każdym
zbiorze zwartym współczynnik po współczynniku
do operatora oscylatora
\[
 H_0=-\sum_{i=1}^{2m}
       \left(\partial_{z_i}
          -\tfrac14\sum_j\Omega_{ij}z^j\right)^2
       -F.
\]
Tutaj \(\Omega=(\Omega_{ij})\) jest macierzą
krzywizny stycznej o wartościach w parzystych
formach, a \(F\) oznacza działanie krzywizny
skręcenia po przejściu do symboli; znak
przy \(F\) odpowiada zapisowi \(e^{-sH_0}\).
Wyrazy zawierające pochodne krzywizny lub wyższe
potęgi \(y\) mają dodatnią potęgę
\(\varepsilon\) i znikają.
Jądro \(K^\varepsilon_s\) jest otrzymane z
\(K_{\varepsilon^2s}\) przez tę samą zmianę
zmiennych oraz czynnik objętości
\(\varepsilon^{2m}\). Po wybraniu składowej
symbolu stopnia \(2m\) oba te czynniki
się znoszą.

Wzór Mehlera dla jednowymiarowego oscylatora,
zastosowany po formalnym rozszczepieniu
antysymetrycznej macierzy \(\Omega\) na bloki,
daje na przekątnej w \(z=0\)
\begin{equation}\label{eq:hodge-mehler}
 K^0_s(0,0)
 =(4\pi s)^{-m}
 \det\nolimits^{1/2}
 \left(\frac{s\Omega/2}{\sinh(s\Omega/2)}\right)
 e^{sF}.
\end{equation}
Rzeczywiście, dla
\(-\partial_z^2+a^2z^2\) jądro ma postać
\[
 \left(\frac{a}{2\pi\sinh(2as)}\right)^{1/2}
 \exp\!\left[
  -\frac{a\bigl((z^2+w^2)\cosh(2as)-2zw\bigr)}
         {2\sinh(2as)}
 \right].
\]
Bezpośrednie różniczkowanie sprawdza równanie
ciepła, a granica \(s\downarrow0\) jest deltą.
Rozłożenie macierzy na bloki oraz formalne
podstawienie za \(a\) jej wartości własnych
daje wyznacznik w~\eqref{eq:hodge-mehler}.
Wszystkie argumenty funkcji hiperbolicznych są
tu formalnymi parzystymi formami; szereg urywa
się w stopniu \(2m\). Każdy blok daje czynnik
\((u/2)/\sinh(u/2)\), więc wyznacznik jest formą
\(\widehat A\).

Pozostaje uzasadnić, że model oblicza granicę
\emph{rzeczywistego} jądra. Po lokalnym
przedłużeniu współczynników poza kulę
różnica przeskalowanych operatorów ma
współczynniki \(O(\varepsilon)\), pomnożone
przez ustalone wielomiany w \(z\).
Tożsamość Duhamela
\[
 e^{-sH_\varepsilon}-e^{-sH_0}
 =-\int_0^s e^{-(s-u)H_\varepsilon}
      (H_\varepsilon-H_0)e^{-uH_0}\,du
\]
i standardowe oszacowania gaussowskie jądra
cieplnego oraz jego pochodnych dają dla
każdego \(\delta>0\) oszacowanie
\[
 |K^\varepsilon_s(z,0)-K^0_s(z,0)|
 \le C_\delta\varepsilon
       s^{-m+1-\delta}e^{-a|z|^2/s}
 \quad (0<s\le1).
\]
W szczególności na przekątnej przy \(s=1\)
prawa strona ma normę \(O(\varepsilon)\).
Właśnie te
oszacowania, a nie gotowa formuła indeksowa,
są przyjętym wejściem analitycznym.
Stosujemy je do jednostajnie eliptycznych
operatorów drugiego rzędu na \(\R^{2m}\)
z jednostajnie ograniczonymi \(g,g^{-1}\),
wyrazem pierwszego rzędu ograniczonym przez
\(C(1+|z|)\) i wyrazem zerowym ograniczonym
przez \(C(1+|z|)^2\); taką rodzinę dają
\(H_\varepsilon\) po lokalnym wycięciu.
Szczegółowy dowód oszacowania dla
współczynników rosnących liniowo i kwadratowo
podaje \href{https://www.maths.ed.ac.uk/~v1ranick/papers/getzler.pdf}
{Getzler, dodatek A}; lokalne wycięcie nie
zmienia granicy, bo wkład dróg opuszczających
kulę maleje wykładniczo jak \(e^{-c/t}\).

Superślad spinorowy usuwa wszystkie symbole
oprócz składowej najwyższego stopnia:
\(\operatorname{str}_S(a)
=(2/i)^m\langle[a]_{2m},\mathrm{vol}_x\rangle\)
w dodatniej orientacji. Widać to na bazowych
iloczynach generatorów Clifforda: dla mniej
niż \(2m\) generatorów odpowiedni ślad jest
zerowy, a iloczyn wszystkich jest skalarną
wielokrotnością chiralii. Po zastosowaniu tej
reguły do~\eqref{eq:hodge-mehler} czynnik
\((4\pi)^{-m}(2/i)^m=(2\pi i)^{-m}\)
wprowadza normalizację krzywizn \(R/(2\pi i)\)
i \(F^W/(2\pi i)\). Pozostaje dokładnie
składowa najwyższego stopnia formy po prawej
stronie~\eqref{eq:hodge-local-index}.
\end{proof}

Łącząc Lemat~\ref{lem:hodge-mckean-singer}
z Lematem~\ref{lem:hodge-local-index} i przechodząc
z granicą pod całką na zwartej \(M\), otrzymujemy
dla skręconego Diraca
\begin{equation}\label{eq:hodge-twisted-dirac-index}
 \operatorname{ind}_{\C}D_W^+
 =\bigl\langle\widehat A(TM)\operatorname{ch}(W),
                 [M]\bigr\rangle.
\end{equation}
Dowód tej równości użył lokalnego modelu i
rachunku Mehlera; nie odwoływał się do ogólnego
twierdzenia o indeksie dowolnego symbolu.

\subsection{Sygnatura i indeks Diraca w wymiarze osiem}
\label{subsec:hodge-indeks-sygnatura}
\index{twierdzenie o sygnaturze Hirzebrucha}
\index{rodzaj A-daszek@\(\widehat A\)}

\begin{twierdzenie}[Potrzebne postacie wzoru na sygnaturę]
\label{thm:hodge-signature-index}
Dla zamkniętej zorientowanej gładkiej
\(4\)-rozmaitości \(X\) oraz \(8\)-rozmaitości \(Y\)
zachodzą odpowiednio
\begin{equation}\label{eq:hodge-signature-polynomials}
 \operatorname{sign}(X)
   =\frac13\langle p_1(TX),[X]\rangle,\qquad
 \operatorname{sign}(Y)
   =\frac1{45}
       \langle7p_2(TY)-p_1(TY)^2,[Y]\rangle.
\end{equation}
\end{twierdzenie}
\begin{proof}
Na zespolonych formach parzystowymiarowej
rozmaitości działa \(d+\delta\). W wymiarze
\(4k\) nadajmy im gradację
\[
 \tau|_{\Omega^p}=i^{\,2k+p(p-1)}*.
\]
Z relacji dla gwiazdy wynika
\(\tau^2=1\) i
\(\tau(d+\delta)=-(d+\delta)\tau\).
Indeks części z dodatniej do ujemnej
podprzestrzeni \(\tau\) jest, dzięki
Lematowi~\ref{lem:hodge-mckean-singer},
śladem \(\tau\) na harmonicznych formach.
Gwiazda paruje stopnie \(p\) i \(4k-p\),
więc ich ślady poza środkiem się znoszą.
Na stopniu \(2k\) mamy \(\tau=*\).
Wzór~\eqref{eq:hodge-signature} oraz jego
ten sam dowód w wymiarze \(4k\) dają
indeks równy \(\operatorname{sign}(M)\).

Lokalnie \(\Lambda^*T^*M\otimes\C
\cong S\otimes S^*\); lewy czynnik jest
modułem Clifforda, a \(d+\delta\) ma symbol
Diraca skręconego przez całe \(S^*\).
Gradacja \(\tau\) jest chiralią lewego
czynnika. Na rozmaitości niespinowej
oba spinorowe czynniki osobno mogą nie
istnieć, lecz ich iloczyn jest globalną
wiązką form. Stosujemy więc lokalny
rachunek z Lematu~\ref{lem:hodge-local-index}.

Niech \(x_1,\ldots,x_{2k}\) będą formalnymi
pierwiastkami Pontriagina, czyli
\(p(TM)=\prod_j(1+x_j^2)\). Wtedy
\[
 \widehat A(TM)=\prod_j
       \frac{x_j/2}{\sinh(x_j/2)},\qquad
 \operatorname{ch}(S^*)=
       \prod_j2\cosh(x_j/2).
\]
Ich iloczyn to
\(2^{2k}\prod_j((x_j/2)/\tanh(x_j/2))\).
Składowa stopnia \(4k\) jest jednorodna
stopnia \(2k\) w zmiennych \(x_j\);
czynnik \(2^{2k}\) znosi podstawienie
\(x_j\mapsto x_j/2\). Otrzymujemy
\[
 \operatorname{sign}(M^{4k})
  =\left\langle\left[
     \prod_j\frac{x_j}{\tanh x_j}
    \right]_{4k},[M]\right\rangle.
\]
Rozwinięcie
\(x/\tanh x=1+x^2/3-x^4/45+\cdots\)
daje w stopniu cztery \(p_1/3\).
W stopniu osiem otrzymujemy
\[
 \frac19\sum_{a<b}x_a^2x_b^2
 -\frac1{45}\sum_ax_a^4
 =\frac{7p_2-p_1^2}{45},
\]
bo \(\sum_ax_a^4=p_1^2-2p_2\).
To dowodzi obu wzorów.
\end{proof}

\begin{twierdzenie}[Całkowitość potrzebnej liczby
\(\widehat A\)]
\label{thm:hodge-dirac-eight}
Jeśli \(Y^8\) jest zamkniętą spinową
rozmaitością, to
\begin{equation}\label{eq:hodge-ahat-eight}
 \widehat A(Y):=
 \frac{7\langle p_1(TY)^2,[Y]\rangle
       -4\langle p_2(TY),[Y]\rangle}{5760}
 =\operatorname{ind}_{\C}D^+_Y\in\Z.
\end{equation}
W konsekwencji
\begin{equation}\label{eq:hodge-896}
 \langle p_1(TY)^2,[Y]\rangle
      -4\operatorname{sign}(Y)
 =896\,\widehat A(Y).
\end{equation}
\end{twierdzenie}
\begin{proof}
W~\eqref{eq:hodge-twisted-dirac-index}
bierzemy trywialną linię \(W\). Dla
formalnych pierwiastków Pontriagina
\[
 \frac{x/2}{\sinh(x/2)}
 =1-\frac{x^2}{24}+\frac{7x^4}{5760}
  +O(x^6).
\]
Zatem składowa stopnia osiem iloczynu
czterech takich czynników wynosi
\[
 \frac7{5760}\sum_jx_j^4
 +\frac1{576}\sum_{i<j}x_i^2x_j^2
 =\frac{7p_1^2-4p_2}{5760}.
\]
Lewa strona~\eqref{eq:hodge-ahat-eight}
jest indeksem, więc jest całkowita.
Eliminując \(p_2\) między
\eqref{eq:hodge-signature-polynomials}
i~\eqref{eq:hodge-ahat-eight}, dostajemy
\(896\widehat A=p_1^2-4\operatorname{sign}\).
\end{proof}

\subsection{Indeks operatora deformacji ASD}
\label{subsec:hodge-indeks-asd}
\index{operator deformacji ASD}

\begin{twierdzenie}[Indeks deformacji ASD]
\label{thm:hodge-asd-index}
Niech \(X^4\) będzie spójną, zamkniętą, zorientowaną
gładką rozmaitością, \(P\to X\) wiązką
\(\mathrm{SU}(2)\), a \(A\) połączeniem na \(P\).
Dla operatora
\[
 D_A=d_A^*\oplus d_A^+:
 \Omega^1(X;\operatorname{ad}P)\longrightarrow
 \Omega^0(X;\operatorname{ad}P)\oplus
 \Omega^{2,+}(X;\operatorname{ad}P)
\]
oraz \(k=\langle c_2(P),[X]\rangle\) mamy
\begin{equation}\label{eq:hodge-asd-index}
 \operatorname{ind}_{\R}D_A
 =-2\langle p_1(\operatorname{ad}P),[X]\rangle
   -\tfrac32\bigl(\chi(X)+\operatorname{sign}(X)\bigr)
 =8k-3(1-b_1(X)+b_2^+(X)).
\end{equation}
\end{twierdzenie}
\begin{proof}
Symbol \(d_A^*\oplus\sqrt2\,d_A^+\)
jest izomorfizmem dla \(\xi\ne0\):
część równoległą do \(\xi\) wykrywa
kontrakcja \(d_A^*\), a na części
prostopadłej \(\sqrt2\,\pi_+(\xi\wedge-)\)
zachowuje normę. Skalowanie drugiej
składowej celu przez \(\sqrt2\)
nie zmienia jądra ani kojądra, a więc
indeksu. Ten znormalizowany symbol
jest symbolem Diraca.

Na kuli z lokalnym podniesieniem spinowym
istnieją identyfikacje
\[
 (\Lambda^1_\C,\,
   \Lambda^0_\C\oplus\Lambda^{2,+}_\C)
 \cong
 (S^-\otimes S^+,\,S^+\otimes S^+).
\]
W drugim czynniku
\(S^+\otimes S^+
\cong\Lambda^0_\C\oplus\Lambda^{2,+}_\C\):
część antysymetryczna jest linią
skalarną, a symetryczna ma wymiar trzy.
Stąd symbol \(D_A\) jest symbolem
\(D^-\) skręconego przez
\(S^+\otimes(\operatorname{ad}P)_\C\).
Składniki niższego rzędu są zwartymi
zaburzeniami między przestrzeniami
Sobolewa \(H^1\to L^2\) na zwartej
rozmaitości, więc nie zmieniają indeksu
Fredholma. Kompleksyfikacja zachowuje
wymiary jąder i kojąder:
\(\operatorname{ind}_{\R}D_A
=\operatorname{ind}_{\C}(D_A)_\C\).
Na niespinowym \(X\) osobne
\(S^\pm\) nie sklejają się, ale oba
iloczyny \(S^-\otimes S^+\) i
\(S^+\otimes S^+\) są wiązkami
stowarzyszonymi z \(\mathrm{SO}(4)\).
Można więc użyć tej identyfikacji
lokalnie w rachunku cieplnym.

W wymiarze cztery piszemy
\(p_1(TX)=x_1^2+x_2^2\) i
\(e(TX)=x_1x_2\). Wagi \(S^+\)
wynoszą \(\pm(x_1+x_2)/2\), stąd
\[
 \operatorname{ch}(S^+)
 =2+\tfrac14(p_1(TX)+2e(TX)),\qquad
 \widehat A(TX)=1-\tfrac1{24}p_1(TX)
\]
do stopnia cztery. Dla rzeczywistej
wiązki rangi trzy
\(\operatorname{ch}((\operatorname{ad}P)_\C)
=3+p_1(\operatorname{ad}P)\)
w tych samych stopniach. Część
stopnia cztery ich iloczynu to
\[
 \left[
 \widehat A(TX)\operatorname{ch}(S^+)
 \operatorname{ch}((\operatorname{ad}P)_\C)
 \right]_4
 =2p_1(\operatorname{ad}P)
  +\tfrac12p_1(TX)+\tfrac32e(TX).
\]
Ponieważ używamy operatora o chiralii
ujemnej, indeks jest \emph{przeciwny}
do całki tej formy.
Twierdzenie~\ref{thm:hodge-signature-index}
i wzór Cherna--Gaussa--Bonneta
z Rozdziału~\ref{ch:chern-weil} dają
\(\langle p_1(TX),[X]\rangle
=3\operatorname{sign}(X)\) oraz
\(\langle e(TX),[X]\rangle=\chi(X)\).

Pozostaje przeliczyć klasę wiązki
adjungowanej. Formalne pierwiastki
wiązki \(E=P\times_{\mathrm{SU}(2)}\C^2\)
to \(u,-u\), więc \(c_2(E)=-u^2\).
Kompleksyfikacja \(\operatorname{ad}P\)
ma wagi \(2u,0,-2u\), skąd
\[
 p_1(\operatorname{ad}P)
 =-c_2((\operatorname{ad}P)_\C)
 =4u^2=-4c_2(P).
\]
Otrzymujemy pierwszy i następnie
drugi wzór w~\eqref{eq:hodge-asd-index},
bo z dualności Poincarégo
\(\chi+\operatorname{sign}
=2(1-b_1+b_2^+)\).
\end{proof}

Wzór~\eqref{eq:hodge-asd-index} liczy
\emph{wymiar oczekiwany}. Gładkość
przestrzeni moduli wymaga osobnej
regularności \(H_A^2=0\), której
sam indeks nie zapewnia.
Konstrukcję moduli i te wejścia
omawia Rozdział~\ref{ch:donaldson}.

\par\smallskip
{\footnotesize\noindent\emph{Źródła dowodu indeksowego:}
\href{https://www.maths.ed.ac.uk/~v1ranick/papers/getzler.pdf}
{E.~Getzler, \emph{A Short Proof of the Local
Atiyah--Singer Index Theorem} (1986)},
\href{https://doi.org/10.1007/BF01425417}
{M.~Atiyah, R.~Bott i V.~K.~Patodi,
\emph{On the Heat Equation and the Index Theorem} (1973)}.
Identyfikacja symbolu deformacji ma źródło w
\href{https://doi.org/10.1098/rspa.1978.0143}
{M.~Atiyah, N.~Hitchin i I.~Singer,
\emph{Self-duality in Four-dimensional Riemannian
Geometry} (1978)}; powyżej przeliczono ją
w konwencji ASD i \(p_1=-4c_2\)
przyjętej w tym skrypcie.\par}

\par\smallskip
{\footnotesize\noindent\emph{Dalsza lektura:}
\href{https://doi.org/10.2307/1969552}
{K.\ Kodaira, \emph{Harmonic Fields in Riemannian Manifolds
(Generalized Potential Theory)} (1949)},
\href{https://doi.org/10.1090/S0002-9904-1946-08647-4}
{S.\ Bochner, \emph{Vector Fields and Ricci Curvature} (1946)}.\par}
